\subsection{Coding standards}

\subsubsection*{Current conventions}

\begin{itemize}
  \item Strict TypeScript, ES2022 target, ESM for application code, and NodeNext with explicit \Code{.js} relative imports in backend TypeScript.
  \item React function components and hooks; no class components or state library beyond context/local state.
  \item Nest decorators at controller/module boundaries; domain validation is implemented with pure functions receiving \Code{unknown}.
  \item Prepared SQLite statements and parameter binding for values; schema/query fragments are static constants.
  \item Spanish user-visible messages and English identifiers/file names, with domain values in English.
  \item Component files use PascalCase; hooks use \Code{use...}; services/modules use kebab or dot suffixes.
\end{itemize}

No consistent repository-wide formatting convention currently exists. Some source is highly compressed onto one line; other source uses conventional multiline formatting. There is no ESLint, Prettier, Stylelint, EditorConfig, or npm lint/format script.

\subsubsection*{Required invariant-preserving practices}

When changing queue mutations, preserve transaction boundaries, server-side validation, active-number uniqueness, request-ID replay semantics, state broadcast, and the rule that only new/recall operations create announcements. Do not treat UI checks as authoritative validation.

When changing contracts, update both \Code{server/models/state.ts} and \Code{src/types.ts} until shared generation is introduced. When changing SPA paths, update both \Code{App.tsx} and the static route list in \Code{bootstrap.ts}.

\subsubsection*{Recommendations}

Adopt one formatter and linter with checked-in configuration; add type-check/lint scripts and CI; share or generate API types; use stable domain error codes; format JSX and CSS for reviewability. These are recommendations, not current standards.
